\section{Related Work}
\label{sec:related}

\paragraph{Vision Transformers and transfer learning.}
ViT~\cite{dosovitskiy2021vit} applies a Transformer~\cite{vaswani2017attention} to a sequence of $16 \times 16$ image patches.
Supervised pretraining on ImageNet-21k~\cite{deng2009imagenet} with strong augmentation and regularization (AugReg~\cite{steiner2022augreg}) gives backbones that transfer well.
We use the AugReg ViT-B/16 checkpoint from timm~\cite{wightman2019timm} that was also fine-tuned on ImageNet-1k.
Models that are better on ImageNet also transfer better, both with a linear probe and with fine-tuning~\cite{kornblith2019better}.
Common fine-tuning recipes use AdamW~\cite{loshchilov2019adamw}, warmup with cosine decay~\cite{loshchilov2017sgdr}, label smoothing~\cite{szegedy2016rethinking}, stochastic depth~\cite{huang2016stochastic} and layer-wise learning-rate decay~\cite{bao2022beit}; our full fine-tuning recipe uses all of them.

\paragraph{Parameter-efficient fine-tuning (PEFT).}
Adapters add a small bottleneck MLP to every Transformer block~\cite{houlsby2019adapter}.
AdaptFormer~\cite{chen2022adaptformer} instead puts a scaled bottleneck branch \emph{next to} the MLP of each ViT block.
LoRA~\cite{hu2022lora} learns a low-rank update $\Delta W = BA$ for frozen weight matrices, and this update can be merged into the weights after training.
VPT~\cite{jia2022vpt} adds learnable prompt tokens to the input, and BitFit~\cite{benzaken2022bitfit} trains only the bias terms.
An older option is to fine-tune only the top layers, whose features are the most task-specific~\cite{yosinski2014transferable}.
We implement LoRA with the PEFT library~\cite{mangrulkar2022peft} and compare it with Adapter, AdaptFormer, BitFit and top-block fine-tuning under one protocol.

\paragraph{Starting the classifier head.}
When a pretrained model is moved to new classes, the old head is replaced by a new one; ViT starts it at zero~\cite{dosovitskiy2021vit}.
LP-FT~\cite{kumar2022lpft} argues that a random head distorts the pretrained features during fine-tuning.
It therefore first trains a linear probe on frozen features and then fine-tunes everything.
Prototype-based classifiers represent each class by the mean feature of its examples, as in prototypical networks~\cite{snell2017prototypical}.
These normalized class means can be copied (\emph{imprinted}) directly into the weights of a cosine classifier~\cite{qi2018imprinted}.
Our prototype head is such an imprinted head with a chosen temperature.
We compare it with a random head and a linear-probe head in a setting with many classes (10,000), few images per class (45) and a short training budget.
